% !TEX root = ../main.tex
\begin{abstracteng}
  \justifying
  The growing mobility of two-wheeled vehicles in university environments raises safety risks, particularly through non-compliance with helmet usage. Existing manual surveillance is hardly scalable and does not provide precise spatio-temporal insights to support policy evaluation. This study designs and evaluates an end-to-end smart surveillance architecture built on real-time video streaming. The architecture combines YOLOv8 for spatial feature extraction, ByteTrack for multi-object tracking, \textit{Intersection over Area} (IoA) filters and \textit{N-Frame} temporal confirmation to suppress false detections, and an event stream that flows through Apache Kafka and Apache Spark Structured Streaming into a PostgreSQL \textit{data mart} that is visualized on an interactive Grafana dashboard. In experimental evaluation, the detection model reaches a \textit{Mean Average Precision} (mAP@50) of 91.4\%; the edge processing pipeline operates stably at 18.9 FPS with 53\,ms inference latency; and the message broker records a 100\% delivery success rate with a 2.65\,ms median latency. Operational testing on 10 hours of field data (07.00--17.00 local time) reduces 11{,}172 raw detections to 113 actual violation events. The dashboard maps a bimodal density distribution, with an anomaly peak recorded at 14:45--15:15 local time that correlates with rainfall and class transition dynamics. The integrated infrastructure shifts reactive surveillance toward a proactive \textit{Decision Support System} that improves the planning of safety patrols in a campus environment.

  \bigskip
  \noindent
  \textbf{Keywords:} YOLOv8, ByteTrack, Stream Processing, Data Mart, Traffic Safety.
\end{abstracteng}
